% ============================================================
% Préambule — PT_ZETA (article FR)
% Modèle spectral canonique des zéros de zêta via la dynamique
% de persistance PT.
% Source : article_PT_zeta_FR.md (consolidation 2026-05-15).
% Adapté du préambule PT_GEOMETRIC_SPECTRAL/preamble.tex.
% ============================================================

\documentclass[11pt,a4paper]{article}

% --- Encodage et langue ---
\usepackage[utf8]{inputenc}
\usepackage[T1]{fontenc}
\usepackage[french]{babel}
\usepackage{csquotes}

% --- Mathématiques ---
\usepackage{amsmath,amssymb,amsthm,mathtools}
\usepackage{bm}
\usepackage{dsfont}

% --- Mise en page ---
\usepackage[margin=2.5cm]{geometry}
\usepackage{microtype}
\usepackage{enumitem}

% --- Tableaux et figures ---
\usepackage{booktabs}
\usepackage{longtable}
\usepackage{array}
\usepackage{graphicx}
\usepackage{tikz}
\usetikzlibrary{arrows.meta,positioning,fit,calc}

% --- Bibliographie ---
\usepackage[backend=biber,style=numeric,sorting=none]{biblatex}
\addbibresource{references.bib}

% --- Hyperliens ---
\usepackage[colorlinks=true,linkcolor=blue!60!black,citecolor=green!50!black,urlcolor=blue!50!black]{hyperref}
\usepackage{cleveref}

% --- Environnements de théorème ---
\theoremstyle{plain}
\newtheorem{theoreme}{Théorème}[section]
\newtheorem{proposition}[theoreme]{Proposition}
\newtheorem{lemme}[theoreme]{Lemme}
\newtheorem{corollaire}[theoreme]{Corollaire}

\theoremstyle{definition}
\newtheorem{definition}[theoreme]{Définition}
\newtheorem{convention}[theoreme]{Convention}

\theoremstyle{remark}
\newtheorem{remarque}[theoreme]{Remarque}
\newtheorem{exemple}[theoreme]{Exemple}
\newtheorem{conjecture}[theoreme]{Conjecture}

% --- Macros PT canoniques ---
\newcommand{\qplus}{q_{+}}
\newcommand{\qminus}{q_{-}}
\newcommand{\mustar}{\mu^{*}}
\newcommand{\sper}{s}

\newcommand{\gammap}[1]{\gamma_{#1}}
\newcommand{\gammathree}{\gamma_{3}}
\newcommand{\gammafive}{\gamma_{5}}
\newcommand{\gammaseven}{\gamma_{7}}

\newcommand{\thetap}[1]{\theta_{#1}}
\newcommand{\sinp}[1]{\sin^{2}\theta_{#1}}

% --- T-series ---
\newcommand{\Tzero}{\textbf{T0}}
\newcommand{\Tone}{\textbf{T1}}
\newcommand{\Ttwo}{\textbf{T2}}
\newcommand{\Tthree}{\textbf{T3}}
\newcommand{\Tfour}{\textbf{T4}}
\newcommand{\Tfive}{\textbf{T5}}
\newcommand{\Tsix}{\textbf{T6}}

% --- BA-series ---
\newcommand{\BAzero}{\textbf{BA0}}
\newcommand{\BAone}{\textbf{BA1}}
\newcommand{\BAtwo}{\textbf{BA2}}
\newcommand{\BAthree}{\textbf{BA3}}
\newcommand{\BAfour}{\textbf{BA4}}
\newcommand{\BAfive}{\textbf{BA5}}

\newcommand{\Lzero}{\textbf{L0}}
\newcommand{\THM}{\textsc{[Thm]}}

% --- Macros spécifiques RH spectral ---
\newcommand{\Rres}{R}                              % facteur résiduel R(s)
\newcommand{\HBK}{H_{\mathrm{BK}}}                 % Berry-Keating xp
\newcommand{\HPTBK}{H_{\mathrm{PT\text{-}BK}}}     % Berry-Keating PT (umellin)
\newcommand{\DRop}{D_R}                            % opérateur diagonal
\newcommand{\kappap}{\kappa_p}                     % amplitude spectrale
\newcommand{\Ap}{A_p}                              % amplitude PT
\newcommand{\Treg}{T}                              % trace régularisée
\newcommand{\NRvM}{N_{\mathrm{RvM}}}               % densité Riemann-von Mangoldt
\newcommand{\NPT}{N_{\mathrm{PT}}}                 % densité PT
\newcommand{\lambdaH}{\lambda_H}                   % auto-couplage Higgs
\newcommand{\DeltaMaslov}{\Delta_N^{\mathrm{Maslov}}}
\newcommand{\Hilb}{\mathcal{H}}
\newcommand{\Primes}{\mathcal{P}}
\newcommand{\Reals}{\mathbb{R}}
\newcommand{\Cplx}{\mathbb{C}}
\newcommand{\Naturals}{\mathbb{N}}
\newcommand{\Integers}{\mathbb{Z}}
\newcommand{\zetaplus}{\zeta_{+}}
\newcommand{\zetaminus}{\zeta_{-}}
\newcommand{\umax}{u_{\max}}
\newcommand{\umin}{u_{\min}}
\newcommand{\eps}{\varepsilon}
\newcommand{\epsstar}{\varepsilon_{*}}
\DeclareMathOperator{\Tr}{Tr}
\DeclareMathOperator{\Reop}{Re}
\DeclareMathOperator{\Imop}{Im}
\DeclareMathOperator{\sgn}{sgn}

% --- Métadonnées article ---
\title{Mod\`ele spectral canonique des z\'eros de $\zeta$ via la dynamique de persistance PT}
\author{Yan Senez\\\texttt{yan.senez@protonmail.com}}
\date{Mai 2026}
